% Figure 16.1 : la structure d'un fichier kubeconfig (celui de minikube, chapitre 16).
\documentclass[tikz,border=4pt]{standalone}
\usepackage{quai}
\begin{document}
\begin{tikzpicture}[x=1cm,y=1cm,
    sec/.style={font=\titrefig\normalsize, anchor=west},
    e/.style={qbox, font=\ttfamily\scriptsize, text width=5.5cm, align=left, inner sep=4pt},
    cl/.style={e, draw=QBleu, fill=QBleuBg},
    us/.style={e, draw=QOcre, fill=QOcreBg},
    cx/.style={e, draw=QSarcelle, fill=QSarcelleBg, text width=3.9cm},
    lien/.style={qlbl, font=\scriptsize, fill=QPaper, inner sep=1pt}]

  \node[sec, text=QBleu] at (-0.2,3.2) {clusters : où ?};
  \node[cl] (c1) at (2.7,2.2) {name: minikube\\server: https://192.168.49.2:8443\\certificate-authority: ca.crt};
  \node[cl] (c2) at (2.7,0.7) {name: kind-listify\\server: https://127.0.0.1:...};

  \node[sec, text=QOcre] at (-0.2,-0.7) {users : qui ?};
  \node[us] (u1) at (2.7,-1.7) {name: minikube\\client-certificate: client.crt\\client-key: client.key};
  \node[us] (u2) at (2.7,-3.2) {name: kind-listify\\...};

  \node[sec, text=QSarcelle] at (7.6,-3.5) {contexts : quel couple ?};
  \node[cx] (x1) at (9.8,0.2) {name: minikube\\cluster: minikube\\user: minikube\\namespace: default};
  \node[cx] (x2) at (9.8,-2.0) {name: kind-listify\\cluster: kind-listify\\user: kind-listify};

  \draw[qarr, draw=QBleu] (x1.west) -- (c1.east);
  \draw[qarr, draw=QOcre] (x1.west) -- (u1.east);
  \draw[qdarr, draw=QMuted] (x2.west) -- (c2.east);
  \draw[qdarr, draw=QMuted] (x2.west) -- (u2.east);

  \node[qbox, font=\ttfamily\footnotesize, draw=QInk, fill=QPaper, inner sep=5pt] (cur) at (9.8,2.4) {current-context: minikube};
  \draw[qarr] (cur) -- node[lien, right] {choisit} (x1.north);
\end{tikzpicture}
\end{document}
